% Section 2.10, from lectures/L02/appendix/b_exercises.md.

\section{Exercises}
\label{c2:sec:exercises}

\begin{aside}
\textbf{How to work these.} Each exercise is tagged with the kinds of work it asks for:
\kinds{Paper} by hand, \kinds{CircuitVerse} in the browser, \kinds{VHDL} in a module of your own.

\textbf{How to check your work.} Every exercise that asks you to write a VHDL module ships a
self-checking testbench under \file{lectures/L02/exercises} in the companion repository. Write your
module in its \file{exercises/<module>} directory, using the entity name and the \textbf{port order}
the exercise specifies, then run it with GHDL; \secref{c2:app:c} has the three commands.

No FPGA board is needed for any exercise. Your job is to get the VHDL right, and the testbench is
how you confirm it. Short answers to the parts you work on paper are in
Appendix~\ref{app:answers}; reference solutions to the VHDL are published in the companion
repository, under \file{lectures/L02/solutions}. Try each exercise yourself before looking at
either.
\end{aside}

% ------------------------------------------------------------------------------------------
\exerciseset{Karnaugh Maps}
\label{c2:set:1}

\exercise{A three-variable map}{Paper, CircuitVerse}
\label{c2:ex:1}
Derive a minimized equation for \code{X} from the Karnaugh map below, then realize the
corresponding gate network and simulate it in CircuitVerse:

\begin{narrowtable}{@{}cc@{}}
  \thead{ABC} & \thead{X} \\
  \midrule
  000 & 1 \\
  001 & 1 \\
  010 & 0 \\
  011 & 0 \\
  100 & 1 \\
  101 & 1 \\
  110 & 0 \\
  111 & 0 \\
\end{narrowtable}

\note[Tip] Lay \code{AB} down the rows in Gray-code order (\code{00, 01, 11, 10}) and \code{C}
across the columns, exactly as in \secref{c2:a1}'s worked example, before you start grouping
\code{1}s.

\exercise{A four-variable map}{Paper, CircuitVerse}
\label{c2:ex:2}
Derive a minimized equation for \code{X} from the Karnaugh map below, then realize the
corresponding gate network and simulate it in CircuitVerse:

\begin{narrowtable}{@{}cc@{}}
  \thead{ABCD} & \thead{X} \\
  \midrule
  0000 & 0 \\
  0001 & 1 \\
  0010 & 0 \\
  0011 & 1 \\
  0100 & 0 \\
  0101 & 0 \\
  0110 & 0 \\
  0111 & 0 \\
  1000 & 0 \\
  1001 & 1 \\
  1010 & 0 \\
  1011 & 1 \\
  1100 & 1 \\
  1101 & 0 \\
  1110 & 1 \\
  1111 & 0 \\
\end{narrowtable}

\note[Tip] With four variables, put two of them (in Gray code) on each axis, so you get a
$4\times4$ grid. Groups can still wrap around every edge of the grid, not just left-right; remember
to check the top/bottom wraparound too.

% ------------------------------------------------------------------------------------------
\needspace{16\baselineskip}
\exerciseset{Circuits in VHDL}
\label{c2:set:2}

\exercise{A three-output network}{Paper, CircuitVerse, VHDL}
\label{c2:ex:3}
A gate network has:
\begin{itemize}
  \item Four inputs, \code{ABCD}.
  \item Three outputs, \code{XYZ}.
\end{itemize}
It is given by the truth table below:

\begin{narrowtable}{@{}cc@{}}
  \thead{ABCD} & \thead{XYZ} \\
  \midrule
  0000 & 001 \\
  0001 & 010 \\
  0010 & 001 \\
  0011 & 011 \\
  0100 & 100 \\
  0101 & 110 \\
  0110 & 100 \\
  0111 & 111 \\
  1000 & 101 \\
  1001 & 110 \\
  1010 & 101 \\
  1011 & 111 \\
  1100 & 000 \\
  1101 & 010 \\
  1110 & 000 \\
  1111 & 011 \\
\end{narrowtable}

\begin{parts}
  \item Treat \code{X}, \code{Y}, and \code{Z} as three separate single-output functions of
    \code{ABCD}, and derive a minimized equation for each via a Karnaugh map.
  \item Realize the resulting gate network and simulate it in CircuitVerse; confirm all three
    outputs against the truth table above.
  \item Implement the design in VHDL, in a module of your own with four \code{std_logic} inputs and
    three \code{std_logic} outputs (\secref{c1:a4} and \secref{c2:a2} show the pattern for turning a
    derived equation into concurrent VHDL assignments).
  \item Run the testbench (see the aside at the start of these exercises) and confirm all 16 rows
    pass.
\end{parts}
If a row fails, the assertion message names the input combination:
\begin{itemize}
  \item Go back to the Karnaugh map for whichever of \code{X}, \code{Y}, \code{Z} is wrong.
  \item A single wrong output on a single row is nearly always one mis-grouped \code{1} in one map.
\end{itemize}
\note[Tip] Two of the three outputs here turn out to depend on only one or two of the four inputs
once minimized. Don't assume every output needs all four.

\modulefigure{xyz_logic}

\selfcheck name your VHDL entity \code{xyz_logic}, with ports \code{a}, \code{b}, \code{c},
\code{d} (in) and \code{x}, \code{y}, \code{z} (out), declared in that order; its testbench is in
\file{exercises/xyz_logic} and checks all 16 rows of the truth table above.

% ------------------------------------------------------------------------------------------
\exerciseset{Multiplexers}
\label{c2:set:3}

\exercise{The 4-to-1 mux by hand}{Paper, CircuitVerse}
\label{c2:ex:4}
A 4-to-1 multiplexer has:
\begin{itemize}
  \item Four data inputs, \code{A}--\code{D}.
  \item Two selector bits, \code{S[1:0]}.
  \item One output, \code{X}.
\end{itemize}
Take the selector to name its data input directly: \code{S[1:0] = "00"} selects \code{A},
\code{"01"} selects \code{B}, \code{"10"} selects \code{C}, and \code{"11"} selects \code{D}. That is
the mapping Exercise~\ref{c2:ex:6} implements.

\begin{parts}
  \item Write out the multiplexer's truth table (four rows, one per selector combination).
  \item Derive the sum-of-products equation for \code{X} in terms of \code{A}--\code{D} and
    \code{S[1:0]} (\secref{c2:a4} works the same way for the 8-to-1 case; halve it here).
  \item Build and simulate the network in CircuitVerse, and confirm \code{X} follows each data input
    in turn as you change \code{S[1:0]}.
\end{parts}
Keep your answer: Exercise~\ref{c2:ex:6} below implements this exact circuit in VHDL, once
\secref{c2:a5} has given you the \code{process} and \code{case} statements it needs.

\exercise{A 2-to-1 mux with \code{if}}{VHDL}
\label{c2:ex:5}
\Secref{c2:a5} writes the 2:1 multiplexer using a \code{case} statement. Write the same circuit
again with an \code{if}/\code{else} instead, and confirm for yourself that the two describe identical
hardware.

Write an entity named \code{mux2} with:

\begin{booktable}{@{}p{4em}p{4.8em}p{5.2em}L@{}}
  \thead{Port} & \thead{Direction} & \thead{Type} & \thead{Description} \\
  \midrule
  \code{d0} & in & \code{std_logic} & Data input. \\
  \code{d1} & in & \code{std_logic} & Data input. \\
  \code{sel} & in & \code{std_logic} & Selector. \\
  \code{x} & out & \code{std_logic} & The selected data input. \\
\end{booktable}

\noindent The selector routes the data input whose number it names:
\begin{itemize}
  \item \code{sel = '0'} selects \code{d0}.
  \item \code{sel = '1'} selects \code{d1}.
  \item Note this is the opposite way round from \secref{c2:a5}'s \code{case} example, which selects
    its \emph{second} input \code{b} on \code{'0'}. Which input a selector value names is a
    convention, not a rule, so read it off the specification every time rather than assuming.
\end{itemize}
Implement the multiplexer using a process-based architecture:
\begin{itemize}
  \item Include \code{d0}, \code{d1}, and \code{sel} in the sensitivity list.
  \item Use an \code{if}/\code{else} statement on \code{sel}.
  \item Assign \code{x} in every branch.
  \item Do not use:
    \begin{itemize}
      \item A conditional concurrent assignment (\code{x <= d1 when sel = '1' else d0;}).
      \item A \code{with...select} statement:
        \begin{itemize}
          \item a concurrent form of \code{case} this book doesn't cover.
          \item mentioned only so you don't reach for it if you've met it elsewhere.
        \end{itemize}
    \end{itemize}
\end{itemize}

\modulefigure{mux2}

\selfcheck name your entity \code{mux2}, with ports \code{d0}, \code{d1}, \code{sel} (in) and
\code{x} (out), declared in that order; its testbench is in \file{exercises/mux2} and checks all
eight combinations of the three inputs.

\exercise{The 4-to-1 mux in VHDL}{VHDL}
\label{c2:ex:6}
Extend Exercise~\ref{c2:ex:5} into a 4:1 multiplexer named \code{mux4}.

This is the circuit you derived, drew, and simulated by hand in Exercise~\ref{c2:ex:4} above.
Re-derive the \code{case} branches yourself for four inputs and two selector bits rather than
copying the 8-to-1 example from \secref{c2:a5}.

The entity must have:

\begin{booktable}{@{}p{4em}p{4.8em}p{14.6em}L@{}}
  \thead{Port} & \thead{Direction} & \thead{Type} & \thead{Description} \\
  \midrule
  \code{d0} & in & \code{std_logic} & Data input. \\
  \code{d1} & in & \code{std_logic} & Data input. \\
  \code{d2} & in & \code{std_logic} & Data input. \\
  \code{d3} & in & \code{std_logic} & Data input. \\
  \code{sel} & in & \code{std_logic_vector(1 downto 0)} & Selector. \\
  \code{x} & out & \code{std_logic} & The selected data input. \\
\end{booktable}

\noindent As in Exercise~\ref{c2:ex:5}, the selector names its data input directly:
\begin{itemize}
  \item \code{"00"} selects \code{d0}.
  \item \code{"01"} selects \code{d1}.
  \item \code{"10"} selects \code{d2}.
  \item \code{"11"} selects \code{d3}.
  \item Again, that is the opposite mapping from \secref{c2:a5}'s 8-to-1 worked example, where
    \code{"000"} routes \code{inputs(7)}. Derive your \code{case} branches from the four lines above,
    not from that listing.
\end{itemize}
Implement the multiplexer using a process:
\begin{itemize}
  \item Include all data inputs and \code{sel} in the sensitivity list.
  \item Use a \code{case} statement on \code{sel}, one branch per selector value above.
  \item Include a \code{when others} branch.
    \begin{itemize}
      \item Choose a safe output value for select inputs containing values such as \code{'X'},
        \code{'U'}, or \code{'Z'}.
    \end{itemize}
  \item Do not replace the \code{case} statement with an \code{if}/\code{elsif} chain.
\end{itemize}

\modulefigure{mux4}

\selfcheck name your entity \code{mux4}, with ports \code{d0}, \code{d1}, \code{d2}, \code{d3},
\code{sel} (\code{std_logic_vector(1 downto 0)}) and \code{x} (out), declared in that order; its
testbench is in \file{exercises/mux4} and sweeps all four selector values against all sixteen
data-input combinations.

% ------------------------------------------------------------------------------------------
\exerciseset{Submodules and 7-Segment Displays}
\label{c2:set:4}
Exercises~\ref{c2:ex:7} and~\ref{c2:ex:8} build one design in two halves: a module that drives a
single 7-segment display, and a top level that uses two copies of it to show a hexadecimal digit on
each. Do them in order: Exercise~\ref{c2:ex:8} instantiates what you write in
Exercise~\ref{c2:ex:7}.

A 7-segment display is seven LEDs, named \code{a} through \code{g}, arranged around a figure-eight.
Light the right subset and you get a digit. This is the layout, and it is the one you need in order
to work out the codes below:

\begin{console}
      aaaa
     f    b
     f    b
      gggg
     e    c
     e    c
      dddd
\end{console}

\noindent So \code{1} lights \code{b} and \code{c} (the two on the right), and \code{7} lights
\code{a}, \code{b} and \code{c}. Comparing the codes given for those two digits is a quick way to
confirm you have the layout the right way round before you start deriving \code{A} to \code{F}.

On the DE0-CV the displays are wired \textbf{active-low}: driving a segment to \code{'0'} lights it,
and \code{'1'} turns it off. That is why the code for \code{8}, which lights every segment, is
\code{"0000000"}, and why a blank display is \code{"1111111"}.

Throughout, \code{hex(6 downto 0)} maps to segments \code{g}, \code{f}, \code{e}, \code{d},
\code{c}, \code{b}, \code{a}: bit 6 is \code{g} and bit 0 is \code{a}.

\exercise{One 7-segment digit}{VHDL, Paper}
\label{c2:ex:7}
Write a module named \code{display} that drives one 7-segment display with a hexadecimal digit.

The entity must have:

\begin{booktable}{@{}p{4em}p{4.8em}p{14.6em}L@{}}
  \thead{Port} & \thead{Direction} & \thead{Type} & \thead{Description} \\
  \midrule
  \code{number} & in & \code{std_logic_vector(3 downto 0)} & The value to show, \code{0000}
    through \code{1111}. \\
  \code{hex} & out & \code{std_logic_vector(6 downto 0)} & The segment code for that value. \\
\end{booktable}

\noindent So feeding \code{number} the value \code{0111} must put the code for \code{7} on
\code{hex}, and feeding it \code{1110} must put the code for \code{E} on \code{hex}.

Implement it using a process:
\begin{itemize}
  \item Put \code{number} in the sensitivity list.
  \item Use a \code{case} statement on \code{number}, one branch per value.
  \item Include a \code{when others} branch, and blank the display in it.
\end{itemize}
You are given the codes for \code{0} to \code{9}. Work out \code{A} to \code{F} yourself, from the
segment layout above:

\begin{vhdlcode}
constant DISPLAY_0  : std_logic_vector(6 downto 0) := "1000000";
constant DISPLAY_1  : std_logic_vector(6 downto 0) := "1111001";
constant DISPLAY_2  : std_logic_vector(6 downto 0) := "0100100";
constant DISPLAY_3  : std_logic_vector(6 downto 0) := "0110000";
constant DISPLAY_4  : std_logic_vector(6 downto 0) := "0011001";
constant DISPLAY_5  : std_logic_vector(6 downto 0) := "0010010";
constant DISPLAY_6  : std_logic_vector(6 downto 0) := "0000010";
constant DISPLAY_7  : std_logic_vector(6 downto 0) := "1111000";
constant DISPLAY_8  : std_logic_vector(6 downto 0) := "0000000";
constant DISPLAY_9  : std_logic_vector(6 downto 0) := "0010000";
-- Todo: add DISPLAY_A through DISPLAY_F here.
constant DISPLAY_OFF: std_logic_vector(6 downto 0) := "1111111";
\end{vhdlcode}

\noindent Two of the six are lower-case on a 7-segment display, because their upper-case forms would
be indistinguishable from a digit: \code{b} would read as \code{8} and \code{d} would read as
\code{0}. \code{A}, \code{C}, \code{E} and \code{F} are upper-case.

All sixteen codes are also sitting in \file{display_tb.vhd}, since a self-checking testbench cannot
check an output without knowing what it should be. So this is not a puzzle you can be locked out of.
Derive the six anyway: it takes a couple of minutes with the segment diagram, and reading them off
the testbench teaches you nothing the \code{case} statement does not. The work of this exercise is
the sixteen-branch \code{case}, not the lookup table.

\modulefigure{display}

\selfcheck name your entity \code{display}, with ports \code{number}
(\code{std_logic_vector(3 downto 0)}) and \code{hex} (out, \code{std_logic_vector(6 downto 0)}),
declared in that order; its testbench is in \file{exercises/display} and checks all sixteen values
against the full code table, naming the digit and both codes when one is wrong.

\exercise{Two digits from submodules}{VHDL}
\label{c2:ex:8}
Write a module named \code{hex_display} that shows a two-digit hexadecimal number on two displays,
using two instances of your \code{display} from Exercise~\ref{c2:ex:7}.

The entity must have:

\begin{booktable}{@{}p{4em}p{4.8em}p{14.6em}L@{}}
  \thead{Port} & \thead{Direction} & \thead{Type} & \thead{Description} \\
  \midrule
  \code{input} & in & \code{std_logic_vector(7 downto 0)} & Eight slide switches: the top four are
    the left digit, the bottom four the right. \\
  \code{hex1} & out & \code{std_logic_vector(6 downto 0)} & Segment code for the left digit. \\
  \code{hex0} & out & \code{std_logic_vector(6 downto 0)} & Segment code for the right digit. \\
\end{booktable}

\noindent When it works:
\begin{itemize}
  \item the number on \code{input(7 downto 4)} appears on \code{hex1}.
  \item the number on \code{input(3 downto 0)} appears on \code{hex0}.
\end{itemize}
Build it with submodules, following \secref{c2:a6}:
\begin{itemize}
  \item Create two instances of \code{display}, labelled \code{display1} and \code{display0}.
  \item Connect \code{input(7 downto 4)} and \code{hex1} to \code{display1}.
  \item Connect \code{input(3 downto 0)} and \code{hex0} to \code{display0}.
  \item Use positional \code{port map}, as everywhere else in this book.
  \item Write no logic of your own in \code{hex_display}. If you find yourself writing a \code{case}
    statement here, you are duplicating Exercise~\ref{c2:ex:7} rather than reusing it.
\end{itemize}
Copy your \file{display.vhd} from Exercise~\ref{c2:ex:7} into the exercise directory, since
\code{hex_display} cannot be analyzed without it. This is the first design in the book built from
more than one file, so name \file{display.vhd} first, ahead of the file that instantiates it:

\begin{shell}
ghdl -a --std=93 display.vhd hex_display.vhd hex_display_tb.vhd
ghdl -e --std=93 hex_display_tb
ghdl -r --std=93 hex_display_tb --assert-level=error --stop-time=10ms
\end{shell}

\noindent\Secref{c2:c5} explains why the order matters.

\modulefigure{hex_display}

\selfcheck name your entity \code{hex_display}, with ports \code{input}
(\code{std_logic_vector(7 downto 0)}), \code{hex1} (out) and \code{hex0} (out, both
\code{std_logic_vector(6 downto 0)}), declared in that order; its testbench is in
\file{exercises/hex_display} and sweeps all 256 input values, checking each half independently so a
design that swaps the two digits is caught rather than passing on the values where both halves
happen to match.

\exercise{A digit as gates}{Paper, CircuitVerse}
\label{c2:ex:9}
\textbf{If you have time.} Realize one digit as a gate network, the way Chapter~\ref{c1:ch} and
\secref{c2:a1} had you do it:
\begin{itemize}
  \item Take the four inputs and one segment, say \code{a}, and derive the Boolean equation for it
    from the 16-row truth table. A Karnaugh map is worth using here: most segments minimize
    substantially.
  \item Repeat for as many of the remaining six segments as you have patience for, and build the
    result in CircuitVerse.
  \item Then copy the whole gate network to drive a second display.
\end{itemize}
That copy is the point of the exercise. Duplicating a block of gates is exactly what instantiating
\code{display} a second time did in Exercise~\ref{c2:ex:8}, and it is worth having done once by hand
to see what the one line of VHDL stands for.

% ------------------------------------------------------------------------------------------
\exerciseset{Naming an Intermediate Result}
\label{c2:set:5}

\exercise{Naming a sub-expression}{VHDL, Paper}
\label{c2:ex:10}
Write an entity named \code{combo_logic} that implements the Boolean function:
\[ x = ab + c'd \]
\Secref{c2:a2} made the case that naming a sub-expression with an internal signal costs nothing and
buys readability. This exercise is where you confirm that rather than take it on trust: you write
the function twice, once as a single expression and once with the \code{c'd} term pulled out into
an internal signal named \code{y}.

A signal declared inside an architecture is a wire. It is not a variable, it does not store
anything, and naming a sub-expression with one does not add a gate, a delay, or a flip-flop to the
circuit. It only gives a piece of the network a name, so that you, reading it later, can see what
that piece is for. Both architectures below synthesize to the same gates, and the provided
testbench passes against either.

The entity must have:

\begin{narrowtable}{@{}lll@{}}
  \thead{Port} & \thead{Direction} & \thead{Type} \\
  \midrule
  \code{a} & in & \code{std_logic} \\
  \code{b} & in & \code{std_logic} \\
  \code{c} & in & \code{std_logic} \\
  \code{d} & in & \code{std_logic} \\
  \code{x} & out & \code{std_logic} \\
\end{narrowtable}

\begin{parts}
  \item Implement the entire function using a single concurrent signal assignment:
    \begin{itemize}
      \item Do not declare any internal signals.
      \item Translate each Boolean operation into its corresponding VHDL operator.
    \end{itemize}
  \item Rewrite the architecture using an internal signal named \code{y}:
    \begin{itemize}
      \item Use \code{y} to represent the term:
        \[ c'd \]
      \item Assign the final expression to \code{x}.
      \item This is the same shape as \code{gate_network} in \secref{c2:a2}, which names its
        \code{AND} term \code{ab}.
    \end{itemize}
\end{parts}
Confirm that the two architectures describe the same circuit.

\note[Tip] Two descriptions represent the same combinational circuit when they produce the same
output for every possible input combination. With four inputs:
\begin{itemize}
  \item There are 16 possible input combinations.
  \item Write the complete truth table.
  \item Verify that both architectures produce the same value of \code{x} in every row.
\end{itemize}

\modulefigure{combo_logic}

\selfcheck name your entity \code{combo_logic}, with ports \code{a}, \code{b}, \code{c}, \code{d}
(in) and \code{x} (out), declared in that order; its testbench is in \file{exercises/combo_logic}
and checks all 16 rows. It passes against either architecture, which is the point of the comparison
above.
